%------------------------------------------------------------------------------%
\addtocounter{exemplo}{1}
\begin{frame}{Exemplo \theexemplo}
  Encontre o ponto da esfera \(x^2+y^2+z^2=4\) mais distante do ponto \((1, -1, 1)\)

  \textit{Dica: não é necessário avaliar completamente a distância para determinar o máximo}
\end{frame}

%------------------------------------------------------------------------------%
\begin{frame}{Exemplo \theexemplo}
  Queremos minimizar a distância
  \[
    d(x, y, z) = \sqrt{(x-1)^2 + (y+1)^2 + (z-1)^2}
  \]
  Como a função raiz quadrada é crescente podemos minimizar a função
  \[
    f(x, y, z) = (x-1)^2 + (y+1)^2 + (z-1)^2
  \]
  suas derivadas parciais são
  \[
    f_x(x, y, z) = 2(x-1)
    \qquad
    f_y(x, y, z) = 2(y+1)
    \qquad
    f_z(x, y, z) = 2(z-1)
  \]
\end{frame}

%------------------------------------------------------------------------------%
\begin{frame}{Exemplo \theexemplo}
  As derivadas parciais da função \(g(x, y, z)=x^2+y^2+z^2\) são
  \[
    g_x(x, y, z) = 2x
    \qquad
    g_y(x, y, z) = 2y
    \qquad
    g_z(x, y, z) = 2z
  \]
\end{frame}

%------------------------------------------------------------------------------%
\begin{frame}{Exemplo \theexemplo}
  Aplicando os Multiplicadores de Lagrange, precisamos resolver o sistema
  \[
  \begin{cases}
    2(x-1) = 2\lambda x \\
    2(y+1) = 2\lambda y \\
    2(z-1) = 2\lambda z \\
    x^2 + y^2 + z^2 = 4
  \end{cases}
  \]
\end{frame}

%------------------------------------------------------------------------------%
\begin{frame}{Exemplo \theexemplo}
  Das três primeiras equações obtemos
  \begin{align*}
    2(x-1)       & = 2\lambda x          & 2(y+1)       & = 2\lambda y          & 2(z-1)       & = 2\lambda z          \\
    x-1          & = \lambda x           & y+1          & = \lambda y           & z-1          & = \lambda z           \\
    x(1-\lambda) & = 1                   & y(1-\lambda) & = -1                  & z(1-\lambda) & = 1                   \\
    x            & = \frac{1}{1-\lambda} & y            & =\frac{-1}{1-\lambda} & z            & = \frac{1}{1-\lambda}
  \end{align*}
\end{frame}

%------------------------------------------------------------------------------%
\begin{frame}{Exemplo \theexemplo}
  Sabemos que \(1-\lambda\neq 0\) pois \(x(1-\lambda) = 1\).

  Substituindo na última equação
  \begin{align*}
    x^2 + y^2 + z^2 & = 4
    \\[2mm]
    \left(\frac{1}{1-\lambda}\right)^2 +
    \left(\frac{-1}{1-\lambda}\right)^2 +
    \left(\frac{1}{1-\lambda}\right)^2 & = 4
    \\[2mm]
    \frac{1 + 1 + 1}{\left(1-\lambda\right)^2} & = 4
    \\[2mm]
    \frac{1}{\left(1-\lambda\right)^2} & = \frac{4}{3}
    \\[2mm]
    \frac{1}{1-\lambda} & = \pm\sqrt{\frac{4}{3}}
    = \pm\frac{2}{\sqrt{3}}
  \end{align*}
\end{frame}

%------------------------------------------------------------------------------%
\begin{frame}{Exemplo \theexemplo}
  Para \(\frac{1}{1-\lambda} = \frac{2}{\sqrt{3}}\)
  \[
    x = \frac{2}{\sqrt{3}}
    \qquad
    y = \frac{-2}{\sqrt{3}}
    \qquad
    z = \frac{2}{\sqrt{3}}
  \]
\end{frame}

%------------------------------------------------------------------------------%
\begin{frame}{Exemplo \theexemplo}
  Para \(\frac{1}{1-\lambda} = -\frac{2}{\sqrt{3}}\)
  \[
    x = \frac{-2}{\sqrt{3}}
    \qquad
    y = \frac{2}{\sqrt{3}}
    \qquad
    z = \frac{-2}{\sqrt{3}}
  \]
\end{frame}

%------------------------------------------------------------------------------%
\begin{frame}{Exemplo \theexemplo}
  Temos dois candidatos
  \(\left( \frac{2}{\sqrt{3}}, \frac{-2}{\sqrt{3}}, \frac{2}{\sqrt{3}}\right)\) e
  \(\left( \frac{-2}{\sqrt{3}}, \frac{2}{\sqrt{3}}, \frac{-2}{\sqrt{3}}\right)\).
\end{frame}

%------------------------------------------------------------------------------%
\begin{frame}{Exemplo \theexemplo}
  Avaliando a função em cada ponto temos
  \begin{align*}
    f\left( \frac{2}{\sqrt{3}}, \frac{-2}{\sqrt{3}}, \frac{2}{\sqrt{3}}\right)
    & = \left(\frac{ 2}{\sqrt{3}}-1\right)^2
    + \left(\frac{-2}{\sqrt{3}}+1\right)^2
    + \left(\frac{ 2}{\sqrt{3}}-1\right)^2 \\[2mm]
    & = \left(\frac{2}{\sqrt{3}}-1\right)^2
    + \left(\frac{2}{\sqrt{3}}-1\right)^2
    + \left(\frac{2}{\sqrt{3}}-1\right)^2 \\[2mm]
    & = 3\left(\frac{2}{\sqrt{3}}-1\right)^2
  \end{align*}
\end{frame}

%------------------------------------------------------------------------------%
\begin{frame}{Exemplo \theexemplo}
  \begin{align*}
    f\left( \frac{-2}{\sqrt{3}}, \frac{2}{\sqrt{3}}, \frac{-2}{\sqrt{3}}\right)
    & = \left(\frac{-2}{\sqrt{3}}-1\right)^2
    + \left(\frac{ 2}{\sqrt{3}}+1\right)^2
    + \left(\frac{-2}{\sqrt{3}}-1\right)^2 \\[2mm]
    & = \left(\frac{2}{\sqrt{3}}+1\right)^2
    + \left(\frac{2}{\sqrt{3}}+1\right)^2
    + \left(\frac{2}{\sqrt{3}}+1\right)^2 \\[2mm]
    & = 3\left(\frac{2}{\sqrt{3}}+1\right)^2
  \end{align*}
\end{frame}

%------------------------------------------------------------------------------%
\begin{frame}{Exemplo \theexemplo}
  Portanto o ponto mais distante é
  \(\left( \frac{-2}{\sqrt{3}}, \frac{2}{\sqrt{3}}, \frac{-2}{\sqrt{3}}\right)\)
  \clearpage
\end{frame}

%------------------------------------------------------------------------------%
